\section{Original coefficients: cancellation and recovery}
\label{stab:coef:section}

Retaining the original crossing coefficients leaves two terms to
control: segments wholly outside the head and segments that cross its
boundary. The toolkit gives two estimates for the latter. Transferring
absolute head masses back to the coefficients yields a square-root
bound; perturbing their pairwise products first yields a fourth-root
bound. These are different estimates for the same coefficient data.

A second choice concerns recovery. A signed mixture directly represents
the displacement, whereas selecting one endpoint from the pooled law
can introduce a factor of two. We first compare these choices for exact
finite laws, then retain approximation errors and a nonzero center tail.
The final subsections convert the resulting estimates into midpoint
bounds.

Unless a statement specifies the forest space $X$, we use either
coordinate model $Z$ and the notation of Section~\ref{stab:toolkit}.
Thus $L(A)$ is the evaluation deficit, $D(A)$ the squared-test slack,
and $O(A),C(A)$ the exterior and crossing parts of a represented atom.
All representations retain their chosen segment families and coefficients.

\subsection{Exact finite laws: two different estimates}
\label{stab:coef:exact}

\begin{proposition}\label{stab:coef:exact-tail}
Let $H$ be finite and ancestral, $x=P_Hx$, $P_Hz=0$, and
$\|x\pm z\|\le1$. Put $d=1-\|x\|\in[0,1]$ and
\[
 \omega_{\mathrm{coef}}(d)=2\bigl(\sqrt{2d}+\sqrt{6d+4\sqrt{2d}}\bigr)
 \qquad(d\ge0).
\]
The original coefficients give both
\begin{align}
 \|z\|&\le(\sqrt{12}+2\sqrt2)\sqrt d,
        \label{stab:coef:square-root}\\
 \|z\|&\le\omega_{\mathrm{coef}}(d).
        \label{stab:coef:fourth-root}
\end{align}
More generally, if $0\le d_0<1$ and $\|x\|\ge1-d_0$, the
right side of \eqref{stab:coef:fourth-root} may be replaced by
$\omega_{\mathrm{coef}}(d_0)$.
\end{proposition}

\begin{proof}
First suppose that $x,z$ are supported in a finite ancestral
$K\supset H$. Choose a head test $p$ with $\rho(p)\le1$ and
$x(p)=\|x\|$; when $x\ne0$ it can be chosen of norm one, and
when $x=0$ take $p=0$. The finite atomic-ball identity represents
each endpoint as a finite law of represented atoms. Let $\tau$ be
the endpoint sign and use the equal mixture of the laws. Then
\begin{equation}\label{stab:coef:exact-laws}
 \mathbb EA=x,\qquad \mathbb E(\tau A)=z,\qquad
 \mathbb E_\pm L=\mathbb E L=d.
\end{equation}
Here the last equality uses $z(p)=0$. The gate coordinates of the
mixed mean vanish. Lemma~\ref{stab:coef:pointwise} therefore gives
\begin{equation}\label{stab:coef:exact-means}
 \beta=\|\mathbb E c\|_2
 =\|(\mathbb E O(A)(e_g))_g\|_2\le\sqrt{2d},
 \qquad \|\mathbb E s\|_2\le2\sqrt{2d}.
\end{equation}
Thus \eqref{stab:coef:absolute-head} and
\eqref{stab:coef:via-head} yield
\begin{equation}\label{stab:coef:twelve}
 \|M_s\|_2^2\le12d,\qquad
 \|M_c\|_2\le\sqrt{12d}+\sqrt{2d}.
\end{equation}
The signed mixture in \eqref{stab:coef:exact-laws}, followed by
Lemma~\ref{stab:coef:tail-assembly}, gives
\[
 \|z\|=\|\mathbb E(\tau Q_HA)\|
 \le\|M_c\|_2+\mathbb E\|O(A)\|
 \le\sqrt{12d}+2\sqrt{2d}.
\]
This signed reconstruction has no conditional factor of two.

For the second route use the original-product estimate
\eqref{stab:coef:via-products}. If
$u_g=\mathbb E(c_g)_+$ and $w_g=\mathbb E(-c_g)_+$, it says
\begin{equation}\label{stab:coef:fourth-moments}
 \sum_g u_gw_g\le d+\sqrt{2d},\quad
 \|u-w\|_2\le\sqrt{2d},\quad
 \|u+w\|_2^2\le6d+4\sqrt{2d}.
\end{equation}
Let $\mathbf1_+$ indicate the plus law in the mixture. Now
$\mathbb E(\mathbf1_+Q_HA)=z/2$. Applying the test-segment proof
of Lemma~\ref{stab:coef:tail-assembly} to
$\mathbf1_+C(A)$, and bounding the exterior term by its pooled
mean, gives
\[
 \frac{\|z\|}{2}
 \le\sqrt{6d+4\sqrt{2d}}+\sqrt{2d}.
\]
The factor of two comes from this indicator reconstruction. This
proves \eqref{stab:coef:fourth-root}. The fourth-root order comes from
the product perturbation, while the leading factor two comes from the
indicator recovery; the two losses have separate causes.

The same loss can be recorded as two separate conditional bounds.
If $d\le d_0$, positivity of the conditional weights and
\eqref{stab:coef:conditional-assembly} give
\begin{equation}\label{stab:coef:fourth-conditional}
 \|\mathbb E_+O(A)\|\le2\sqrt{2d_0},\qquad
 \|\mathbb E_+C(A)\|
 \le2\sqrt{6d_0+4\sqrt{2d_0}}.
\end{equation}
Indeed each conditional mean of a nonnegative quantity is at most
twice its pooled mean. Since $z=\mathbb E_+Q_HA$, adding these
bounds gives $\omega_{\mathrm{coef}}(d_0)$ directly, with the exterior
and crossing errors still distinguished.

For a general $z$, let $K$ increase through finite ancestral sets
containing $H$. Contractivity gives
$\|x\pm P_Kz\|\le1$, and $P_HP_Kz=0$.
Apply the finite calculation and use $P_Kz\to z$ in $Z$.
The function $\omega_{\mathrm{coef}}$ is nondecreasing, so $d\le d_0$
proves the last assertion. This also handles nonattainment in the completed
space without a representation of its endpoints.
\end{proof}

\subsubsection{Signed recovery in the forest}
\label{mola:actual-section}
The product estimate also permits signed recovery. The next forest
statement uses the same absolute-mass budget as
\eqref{stab:coef:fourth-moments}, but retains the endpoint sign instead
of selecting one conditional law. Its function $g$ is exactly
$\omega_{\mathrm{coef}}/2$.

\begin{proposition}\label{mola:actual}
In the forest space $X$, suppose $H$ is finite and ancestor-closed,
$x=P_Hx$, $y=Q_Hy$, and
$\|x+y\|,\|x-y\|\le1$.  Then
\[
 \|y\|\le g(1-\|x\|),\qquad
 g(s)=\sqrt{2s}+\sqrt{6s+4\sqrt{2s}}\quad(s\ge0).
\]
\end{proposition}
\begin{proof}
For finite $y$, use the exact laws of
\eqref{stab:coef:exact-laws}, with $z=y$ and $d=s=1-\|x\|$.
The original crossing coefficients obey
\eqref{stab:coef:fourth-moments}, so
\[
 \|M_c\|_2\le\sqrt{6s+4\sqrt{2s}},\qquad
 \mathbb E\|O(A)\|\le\sqrt{2s}.
\]
The signed reconstruction $y=\mathbb E(\tau Q_HA)$ and
Lemma~\ref{stab:coef:tail-assembly} give
\[
 \|y\|\le\|M_c\|_2+\mathbb E\|O(A)\|\le g(s).
\]
There is no conditional factor of two in this calculation. It keeps
the original-product transfer error, which is responsible for the
fourth-root order. For general $y$, apply this finite estimate to
$P_Ky$ along a finite ancestral exhaustion containing $H$.
Contractivity preserves both endpoint bounds, and norm convergence
proves the assertion. When $x=0$, the head test used in the exact-law
calculation is $p=0$, as before.
\end{proof}

We now return to either coordinate model $Z$.

\begin{corollary}\label{stab:coef:unit-head-lens}
If $x=P_Hx$ has norm one, $P_Hy=0$, and $\|x\pm y\|\le R$
for $R\ge1$, then
\begin{equation}\label{stab:coef:unit-head-lens-eq}
 \|y\|\le(\sqrt{12}+2\sqrt2)\sqrt{R(R-1)}.
\end{equation}
\end{corollary}

\begin{proof}
Apply \eqref{stab:coef:square-root} to $x/R,y/R$, whose deficit
is $h=1-1/R$, and multiply by $R$. In particular, both the mean
coefficient error squared and the mean head slack are at most $2h$;
the wholly exterior term has mean norm at most $\sqrt{2h}$.
The crossing means have norms at most $\sqrt{2h}$ and
$2\sqrt{2h}$ before absolute values, and the unsigned head mean
has squared norm at most $12h$. These are the terms retained in
the stated constant.
\end{proof}

\subsection{Approximate laws in the completed atomic ball}
\label{stab:coef:approximate}

The completed atomic-ball description also lets us approximate the
endpoints themselves, rather than first projecting the displacement
to finite support. The norm errors then remain in the cancellation
equation. The gauge and coordinate-span presentations agree, so both
routes are available in either model. Direct approximation does not
require a completed endpoint to be an exact finite convex combination.

\begin{proposition}\label{stab:coef:approximate-tail}
Let $H$ be finite and ancestral, $x=P_Hx$ have norm one, and
$P_Hy=0$. For $d>0$, the assumptions $\|x\pm y\|\le1+d$ imply
\begin{equation}\label{stab:coef:approximate-tail-eq}
 \|y\|\le(1+d)(7\sqrt2+4)\sqrt d.
\end{equation}
In particular, $\|y\|\le32\sqrt d$ when $0<d\le1$.
\end{proposition}

\begin{proof}
Choose a norming head test $p$ as above. Equivalently, in the gauge
presentation choose a norm-one functional norming $x$ and compose
it with $P_H$; its coordinate values define $p$, since optimizing
over disjoint-segment atoms gives $\sum_i p(S_i)^2\le1$.
For $\epsilon>0$, choose finite convex combinations $w^\pm$ of
atoms with
\[
 \left\|w^\pm-\frac{x\pm y}{1+d}\right\|<\epsilon.
\]
Such approximants follow from the closed atomic-ball description,
not an exact convex representation in the completion. This closure
also follows directly from the strict-gauge property: approximate a
completed unit-ball vector $u$ by finite $u_n$ with
$\|u_n-u\|<1/n$. Then $u_n/(1+2/n)$ has gauge strictly below
one and belongs to the convex hull of atoms, while still converging
to $u$. Use the two endpoint approximants' finite
laws and their equal mixture, and put $k=d+\epsilon$. Since
$y(p)=0$,
\begin{equation}\label{stab:coef:approximate-deficit}
 \mathbb E_\pm L\le1-\frac1{1+d}+\epsilon\le k,
 \qquad \mathbb E L\le k,
 \qquad \left\|\mathbb EA-\frac{x}{1+d}\right\|<\epsilon.
\end{equation}
The gate restriction is contractive and vanishes on $x$. Thus
\eqref{stab:coef:gate-coordinates} yields
\begin{equation}\label{stab:coef:approximate-means}
 \|\mathbb Ec\|_2\le\epsilon+\sqrt{2k},\qquad
 \|\mathbb Es\|_2\le\epsilon+2\sqrt{2k}.
\end{equation}
Use the splitting proof of Lemma~\ref{stab:coef:head-conflict}.
Independence gives the same conflict budget as in
Lemma~\ref{stab:coef:gate-bound}, so
\begin{align}
 \|M_s\|_2^2&\le(\epsilon+2\sqrt{2k})^2+4k,
       \label{stab:coef:approximate-absolute-head}\\
 \|M_c\|_2&\le\epsilon+3\sqrt{2k}+2\sqrt k.
       \label{stab:coef:approximate-absolute-coef}
\end{align}
For the second line, first bound the square root in the first line
by $\epsilon+2\sqrt{2k}+2\sqrt k$, and then use the mean
coefficient error $\sqrt{2k}$.

The exterior term under the plus law has norm at most $\sqrt{2k}$
by \eqref{stab:coef:approximate-deficit}. Apply the convex-product
proof of Lemma~\ref{stab:coef:tail-assembly} to its crossing part.
The conditional absolute coefficient mean is bounded coordinatewise
by twice the pooled mean. Hence
\begin{equation}\label{stab:coef:approximate-conditional}
 \|Q_Hw^+\|\le\sqrt{2k}
       +2(\epsilon+3\sqrt{2k}+2\sqrt k).
\end{equation}
Since $Q_H$ is contractive and $Q_H(x+y)=y$, add $\epsilon$ to
this bound for $\|y\|/(1+d)$. Let $\epsilon\downarrow0$.
This proves \eqref{stab:coef:approximate-tail-eq}. Only when
$d\le1$ may one use $1+d\le2$ to obtain
$2(7\sqrt2+4)\sqrt d<32\sqrt d$.
\end{proof}

